% GENERATED FILE. Edit canonical lessons, then run npm run generate:books.
% canonical-source-hash: fnv1a64:adc7042edd5f7287
% canonical-lessons: TA-C127-natcattiram, TA-C127-pani, TA-C127-iti, TA-C127-minnal, TA-C127-puyal

\chapter{Star, Dew, Thunder, Lightning, Storm}
\label{ch:ta-star-dew-thunder-lightning-storm}
\hlchaptermodality{127}

\begin{chapteropening}
\textbf{By the end of this chapter:} \emph{I can say a star, dew, thunder, lightning and a storm.}

Complete the last lesson of chapter 127: i can say a star, dew, thunder, lightning and a storm.
\end{chapteropening}

\section[\texorpdfstring{\ta{நட்சத்திரம்} (naṭcattiram)}{நட்சத்திரம் (naṭcattiram)}]{\ta{நட்சத்திரம்} (naṭcattiram) — a star}

\label{lesson:TA-C127-natcattiram}



\begin{quote}
Before the new one: say the Tamil for a leaf, then the Tamil for soil.

Say the day words, \textbf{\ta{நாள்}} and the formal \textbf{\ta{தினம்}}, then the night words, \textbf{\ta{இரவு}}, \textbf{\ta{ராத்திரி}}, and \textbf{\ta{இருள்}}, darkness.
\end{quote}

\subsection*{You'll want to know: \ta{நட்சத்திரம்}}
\textbf{\ta{நட்சத்திரம்}} — \emph{naṭcattiram} — \textquotedblleft{}a star\textquotedblright{}.

\textbf{In the family, and next door}

\noindent
\begin{tabularx}{\linewidth}{@{}>{\raggedright\arraybackslash}X>{\raggedright\arraybackslash}X>{\raggedright\arraybackslash}X@{}}
\toprule
\textbf{} & \textbf{word} & \textbf{same root} \\
\midrule
Kannada & \ka{ಚುಕ್ಕಿ} (\emph{cukki}) &  \\
Telugu & \te{చుక్క} (\emph{cukka}) &  \\
Malayalam & \ml{നക്ഷത്രം} (\emph{nakṣatraṁ}) & yes \\
Hindi (neighbour) & \dv{तारा} (\emph{tārā}) &  \\
English & a star &  \\
\bottomrule
\end{tabularx}

\begin{grammarlens}[title={What you've built}]
One more word for this chapter.
\end{grammarlens}

\subsection*{Your turn}
\begin{itemize}
\raggedright
  \item \emph{Say it:} \emph{naṭcattiram}
  \item \emph{Say it:} \emph{naṭcattiram}, once more
  \item \emph{Say it:} say \emph{naṭcattiram}
  \item \emph{Recall:} say the Tamil for a leaf, then the Tamil for soil, then say \emph{naṭcattiram} again
\end{itemize}

\begin{tcolorbox}[breakable,colback=teal!4,colframe=teal!35!black,title={Before you move on}]
What does \ta{நட்சத்திரம்} mean? (\textquotedblleft{}A star\textquotedblright{}.) Say it once more.
\end{tcolorbox}
\section[\texorpdfstring{\ta{பனி} (paṉi)}{பனி (paṉi)}]{\ta{பனி} (paṉi) — dew, frost, snow}

\label{lesson:TA-C127-pani}



\begin{quote}
Before the new one: say the Tamil for soil, then the Tamil for a star.
\end{quote}

\subsection*{You'll want to know: \ta{பனி}}
\textbf{\ta{பனி}} — \emph{paṉi} — \textquotedblleft{}dew, frost, snow\textquotedblright{}.

\textbf{In the family, and next door}

\noindent
\begin{tabularx}{\linewidth}{@{}>{\raggedright\arraybackslash}X>{\raggedright\arraybackslash}X@{}}
\toprule
\textbf{} & \textbf{word} \\
\midrule
Hindi (neighbour) & \dv{ओस} (\emph{os}) \\
English & dew, frost, snow \\
\bottomrule
\end{tabularx}

\begin{grammarlens}[title={What you've built}]
One more word for this chapter.
\end{grammarlens}

\subsection*{Your turn}
\begin{itemize}
\raggedright
  \item \emph{Say it:} \emph{paṉi}
  \item \emph{Say it:} \emph{paṉi}, once more
  \item \emph{Say it:} say \emph{paṉi}
  \item \emph{Recall:} say the Tamil for soil, then the Tamil for a star, then say \emph{paṉi} again
\end{itemize}

\begin{tcolorbox}[breakable,colback=teal!4,colframe=teal!35!black,title={Before you move on}]
What does \ta{பனி} mean? (\textquotedblleft{}Dew, frost, snow\textquotedblright{}.) Say it once more.
\end{tcolorbox}
\section[\texorpdfstring{\ta{இடி} (iṭi)}{இடி (iṭi)}]{\ta{இடி} (iṭi) — thunder}

\label{lesson:TA-C127-iti}



\begin{quote}
Before the new one: say the Tamil for a star, then the Tamil for dew.
\end{quote}

\subsection*{You'll want to know: \ta{இடி}}
\textbf{\ta{இடி}} — \emph{iṭi} — \textquotedblleft{}thunder\textquotedblright{}.

\textbf{In the family, and next door}

\noindent
\begin{tabularx}{\linewidth}{@{}>{\raggedright\arraybackslash}X>{\raggedright\arraybackslash}X@{}}
\toprule
\textbf{} & \textbf{word} \\
\midrule
Kannada & \ka{ಗುಡುಗು} (\emph{guḍugu}) \\
Telugu & \te{ఉరుము} (\emph{urumu}) \\
English & thunder \\
\bottomrule
\end{tabularx}

\begin{grammarlens}[title={What you've built}]
One more word for this chapter.
\end{grammarlens}

\subsection*{Your turn}
\begin{itemize}
\raggedright
  \item \emph{Say it:} \emph{iṭi}
  \item \emph{Say it:} \emph{iṭi}, once more
  \item \emph{Say it:} say \emph{iṭi}
  \item \emph{Recall:} say the Tamil for a star, then the Tamil for dew, then say \emph{iṭi} again
\end{itemize}

\begin{tcolorbox}[breakable,colback=teal!4,colframe=teal!35!black,title={Before you move on}]
What does \ta{இடி} mean? (\textquotedblleft{}Thunder\textquotedblright{}.) Say it once more.
\end{tcolorbox}
\section[\texorpdfstring{\ta{மின்னல்} (miṉṉal)}{மின்னல் (miṉṉal)}]{\ta{மின்னல்} (miṉṉal) — lightning}

\label{lesson:TA-C127-minnal}



\begin{quote}
Before the new one: say the Tamil for dew, then the Tamil for thunder.
\end{quote}

\subsection*{You'll want to know: \ta{மின்னல்}}
\textbf{\ta{மின்னல்}} — \emph{miṉṉal} — \textquotedblleft{}lightning\textquotedblright{}.

\textbf{In the family, and next door}

\noindent
\begin{tabularx}{\linewidth}{@{}>{\raggedright\arraybackslash}X>{\raggedright\arraybackslash}X@{}}
\toprule
\textbf{} & \textbf{word} \\
\midrule
Kannada & \ka{ಮಿಂಚು} (\emph{miṁcu}) \\
Telugu & \te{మెరుపు} (\emph{merupu}) \\
English & lightning \\
\bottomrule
\end{tabularx}

\begin{grammarlens}[title={What you've built}]
One more word for this chapter.
\end{grammarlens}

\subsection*{Your turn}
\begin{itemize}
\raggedright
  \item \emph{Say it:} \emph{miṉṉal}
  \item \emph{Say it:} \emph{miṉṉal}, once more
  \item \emph{Say it:} say \emph{miṉṉal}
  \item \emph{Recall:} say the Tamil for dew, then the Tamil for thunder, then say \emph{miṉṉal} again
\end{itemize}

\begin{tcolorbox}[breakable,colback=teal!4,colframe=teal!35!black,title={Before you move on}]
What does \ta{மின்னல்} mean? (\textquotedblleft{}Lightning\textquotedblright{}.) Say it once more.
\end{tcolorbox}
\section[\texorpdfstring{\ta{புயல்} (puyal)}{புயல் (puyal)}]{\ta{புயல்} (puyal) — a storm}

\label{lesson:TA-C127-puyal}



\begin{quote}
Before the new one: say the Tamil for thunder, then the Tamil for lightning.

Say \textbf{\ta{வானிலை}}, then \textbf{\ta{மழை} \ta{பெய்கிறது}}. Name the rainy season, \textbf{\ta{மழைக்} \ta{காலம்}}, the one that shapes the year here.
\end{quote}

\subsection*{You'll want to know: \ta{புயல்}}
\textbf{\ta{புயல்}} — \emph{puyal} — \textquotedblleft{}a storm\textquotedblright{}.

\textbf{In the family, and next door}

\noindent
\begin{tabularx}{\linewidth}{@{}>{\raggedright\arraybackslash}X>{\raggedright\arraybackslash}X@{}}
\toprule
\textbf{} & \textbf{word} \\
\midrule
Kannada & \ka{ಬಿರುಗಾಳಿ} (\emph{birugāḷi}) \\
Telugu & \te{తుపాను} (\emph{tupānu}) \\
Hindi (neighbour) & \dv{आँधी} (\emph{ā̃dhī}) \\
English & a storm \\
\bottomrule
\end{tabularx}

\begin{grammarlens}[title={What you've built}]
One more word for this chapter.
\end{grammarlens}

\subsection*{Your turn}
\begin{itemize}
\raggedright
  \item \emph{Say it:} \emph{puyal}
  \item \emph{Say it:} \emph{puyal}, once more
  \item \emph{Say it:} say \emph{puyal}
  \item \emph{Recall:} say the Tamil for thunder, then the Tamil for lightning, then say \emph{puyal} again
\end{itemize}

\begin{tcolorbox}[breakable,colback=teal!4,colframe=teal!35!black,title={Before you move on}]
What does \ta{புயல்} mean? (\textquotedblleft{}A storm\textquotedblright{}.) Say it once more.
\end{tcolorbox}
